\documentclass[aspectratio=169, 14pt]{beamer}
\usepackage{stampfly_slides}

\lesson{2}{コントローラ入力}{Controller Input}

\begin{document}

% ------------------------------------------------------------------
\lessondivider{2}{コントローラ入力}{Controller Input}

% ------------------------------------------------------------------
\begin{frame}{今日のゴール / Today's Goal}
  \begin{block}{目標: スティック値を変数に読み取り、演算でモータを個別制御する}
    \begin{itemize}
      \item ESP-NOW 無線通信の仕組みを理解
      \item 変数と四則演算でスティック → モータ Duty を計算
      \item オープンループ手動操縦の限界を体感
    \end{itemize}
  \end{block}
\end{frame}

% ------------------------------------------------------------------
\begin{frame}{ESP-NOW 通信フロー / Communication Flow}
  \begin{center}
    \resizebox{0.8\textwidth}{!}{\documentclass[border=10pt]{standalone}
\usepackage{tikz}
\usetikzlibrary{arrows.meta, positioning, fit, decorations.pathmorphing}

\begin{document}

\definecolor{ctrlcolor}{RGB}{0, 120, 200}
\definecolor{espcolor}{RGB}{120, 80, 180}
\definecolor{apicolor}{RGB}{50, 140, 50}
\definecolor{gray1}{RGB}{120, 120, 120}
\begin{tikzpicture}[
    box/.style={
      rectangle, rounded corners=4pt, draw=#1, line width=1.5pt,
      fill=#1!8, minimum width=40mm, minimum height=18mm,
      font=\large\bfseries, text=#1, align=center
    },
    smallbox/.style={
      rectangle, rounded corners=2pt, draw=#1, line width=1pt,
      fill=#1!6, minimum width=32mm, minimum height=14mm,
      font=\bfseries, text=#1, align=center
    },
    arr/.style={-{Stealth[length=3mm]}, line width=1.5pt, #1},
    lbl/.style={font=\small, text=gray1, align=center},
  ]

  % Controller side
  \node[box=ctrlcolor] (ctrl) at (0, 0) {Controller};

  % Sticks (2x2 grid with axis labels inside)
  \coordinate (stickcenter) at ([yshift=-48mm]ctrl.south);
  \node[smallbox=ctrlcolor, left=6mm of stickcenter, yshift=16mm]
    (s_thr) {Throttle\\[-1pt]{\scriptsize\color{gray1} Left Stick Y}};
  \node[smallbox=ctrlcolor, right=6mm of stickcenter, yshift=16mm]
    (s_roll) {Roll\\[-1pt]{\scriptsize\color{gray1} Right Stick X}};
  \node[smallbox=ctrlcolor, left=6mm of stickcenter, yshift=-16mm]
    (s_yaw) {Yaw\\[-1pt]{\scriptsize\color{gray1} Left Stick X}};
  \node[smallbox=ctrlcolor, right=6mm of stickcenter, yshift=-16mm]
    (s_pitch) {Pitch\\[-1pt]{\scriptsize\color{gray1} Right Stick Y}};

  % Arrows: all sticks -> Controller via single tap point
  \coordinate (ctap) at ([yshift=-12mm]ctrl.south);
  \draw[line width=1.5pt, ctrlcolor] (ctrl.south) -- (ctap);
  \fill[ctrlcolor] (ctap) circle (2.5pt);
  % Upper row: direct vertical to tap horizontal line
  \draw[arr=ctrlcolor] (s_thr.north) -- (s_thr.north |- ctap) -- (ctap);
  \draw[arr=ctrlcolor] (s_roll.north) -- (s_roll.north |- ctap) -- (ctap);
  % Lower row: vertical up then horizontal to tap
  \draw[arr=ctrlcolor] (s_yaw.north) |- (ctap);
  \draw[arr=ctrlcolor] (s_pitch.north) |- (ctap);

  % ESP-NOW wireless
  \node[box=espcolor, right=30mm of ctrl] (espnow) {ESP-NOW\\2.4\,GHz};
  \draw[arr=espcolor, decorate, decoration={snake, amplitude=1.5mm, segment length=6mm}]
    (ctrl) -- node[above=3mm, font=\small\bfseries, text=espcolor] {Wireless} (espnow);

  % StampFly
  \node[box=apicolor, right=30mm of espnow] (stampfly) {StampFly};
  \draw[arr=apicolor] (espnow) -- (stampfly);

  % API outputs (2x2 grid with range inside)
  \coordinate (apicenter) at ([yshift=-48mm]stampfly.south);
  \node[smallbox=apicolor, left=6mm of apicenter, yshift=16mm]
    (api1) {rc\_throttle()\\[-1pt]{\scriptsize\color{gray1} 0.0\,--\,1.0}};
  \node[smallbox=apicolor, right=6mm of apicenter, yshift=16mm]
    (api2) {rc\_roll()\\[-1pt]{\scriptsize\color{gray1} -1.0\,--\,+1.0}};
  \node[smallbox=apicolor, left=6mm of apicenter, yshift=-16mm]
    (api3) {rc\_pitch()\\[-1pt]{\scriptsize\color{gray1} -1.0\,--\,+1.0}};
  \node[smallbox=apicolor, right=6mm of apicenter, yshift=-16mm]
    (api4) {rc\_yaw()\\[-1pt]{\scriptsize\color{gray1} -1.0\,--\,+1.0}};

  % Arrows: StampFly -> all APIs via single tap point
  \coordinate (atap) at ([yshift=-12mm]stampfly.south);
  \draw[line width=1.5pt, apicolor] (stampfly.south) -- (atap);
  \fill[apicolor] (atap) circle (2.5pt);
  \draw[arr=apicolor] (atap) -- (atap |- api1.north) -- (api1.north);
  \draw[arr=apicolor] (atap) -- (atap |- api2.north) -- (api2.north);
  \draw[arr=apicolor] (atap) -| (api3.north);
  \draw[arr=apicolor] (atap) -| (api4.north);

\end{tikzpicture}
\end{document}
}
  \end{center}
\end{frame}

% ------------------------------------------------------------------
\begin{frame}{コントローラ API}
  \begin{table}
    \centering
    \begin{tabular}{lll}
      \hline
      \textbf{関数} & \textbf{説明} & \textbf{値域} \\
      \hline
      \api{rc\_throttle()} & スロットル & 0.0 -- 1.0 \\
      \api{rc\_roll()}     & ロール     & -1.0 -- +1.0 \\
      \api{rc\_pitch()}    & ピッチ     & -1.0 -- +1.0 \\
      \api{rc\_yaw()}      & ヨー       & -1.0 -- +1.0 \\
      \api{is\_armed()}    & ARM状態    & true / false \\
      \api{motor\_set\_duty(id, v)} & モータ個別制御 & id=1--4, v=0.0--1.0 \\
      \hline
    \end{tabular}
  \end{table}
\end{frame}

% ------------------------------------------------------------------
\begin{frame}{ペアリング手順 / Pairing}
  \begin{table}
    \centering
    \begin{tabular}{cp{120mm}}
      \hline
      \textbf{Step} & \textbf{操作} \\
      \hline
      1 & コントローラの \textbf{M5ボタンを押しながら電源ON} → LCD に ``Pairing mode...'' 表示 \\
      2 & StampFly のボタンを \textbf{3秒長押し} → 青LED高速点滅 + ビープ音 \\
      3 & ペアリング完了 → ビープ音が鳴り、次回から自動接続 \\
      \hline
    \end{tabular}
  \end{table}

  \begin{alertblock}{注意}
    \begin{itemize}
      \item 教室では\textbf{一組ずつ}ペアリングする（ブロードキャスト通信のため近くの機体と干渉する可能性）
      \item うまくいかない場合: 両方を再起動して Step 1 からやり直す
    \end{itemize}
  \end{alertblock}
\end{frame}

% ------------------------------------------------------------------
\begin{frame}{モータ配置とミキシング / Motor Layout \& Mixing}
  \begin{columns}[T]
    \begin{column}{0.45\textwidth}
      \centering
      \resizebox{0.95\columnwidth}{!}{\documentclass[border=10pt]{standalone}
\usepackage{tikz}
\usetikzlibrary{arrows.meta, positioning}

\begin{document}

\definecolor{cwcolor}{RGB}{0, 120, 200}    % Blue for CW
\definecolor{ccwcolor}{RGB}{220, 60, 20}   % Red for CCW
\definecolor{bodycolor}{RGB}{80, 80, 80}   % Dark gray
\begin{tikzpicture}[
    motor/.style={circle, draw, line width=1.5pt, minimum size=18mm, font=\large\bfseries},
    cw/.style={motor, draw=cwcolor, text=cwcolor},
    ccw/.style={motor, draw=ccwcolor, text=ccwcolor},
    label/.style={font=\small, text=bodycolor},
    arrow/.style={-{Stealth[length=3mm]}, line width=1.2pt},
  ]

  % Body (X-frame)
  \draw[bodycolor, line width=2pt] (-2.5, 2.5) -- (2.5, -2.5);
  \draw[bodycolor, line width=2pt] (2.5, 2.5) -- (-2.5, -2.5);

  % Center body
  \fill[bodycolor, rounded corners=3pt] (-0.6, -0.6) rectangle (0.6, 0.6);
  \node[font=\footnotesize, text=white] at (0, 0) {ESP32};

  % Front direction arrow
  \draw[arrow, bodycolor] (0, 1.0) -- (0, 2.0);
  \node[label, above] at (0, 2.0) {Front};

  % M1 - Front Right (CCW)
  \node[ccw] (m1) at (2.5, 2.5) {M1};
  \node[label, above right=-1mm] at (m1.north east) {FR};
  \draw[arrow, ccwcolor] ([shift={(-45:1.3)}]m1.center)
    arc[start angle=-45, end angle=225, radius=1.3cm];
  \node[font=\scriptsize, ccwcolor] at (4.2, 2.5) {CCW};

  % M2 - Rear Right (CW)
  \node[cw] (m2) at (2.5, -2.5) {M2};
  \node[label, below right=-1mm] at (m2.south east) {RR};
  \draw[arrow, cwcolor] ([shift={(45:1.3)}]m2.center)
    arc[start angle=45, end angle=-225, radius=1.3cm];
  \node[font=\scriptsize, cwcolor] at (4.2, -2.5) {CW};

  % M3 - Rear Left (CCW)
  \node[ccw] (m3) at (-2.5, -2.5) {M3};
  \node[label, below left=-1mm] at (m3.south west) {RL};
  \draw[arrow, ccwcolor] ([shift={(225:1.3)}]m3.center)
    arc[start angle=225, end angle=495, radius=1.3cm];
  \node[font=\scriptsize, ccwcolor] at (-4.2, -2.5) {CCW};

  % M4 - Front Left (CW)
  \node[cw] (m4) at (-2.5, 2.5) {M4};
  \node[label, above left=-1mm] at (m4.north west) {FL};
  \draw[arrow, cwcolor] ([shift={(135:1.3)}]m4.center)
    arc[start angle=135, end angle=-135, radius=1.3cm];
  \node[font=\scriptsize, cwcolor] at (-4.2, 2.5) {CW};

  % Legend
  \node[anchor=west, font=\small] at (-4.5, -4.5) {%
    \textcolor{cwcolor}{\rule{8pt}{8pt}} CW (clockwise) \quad
    \textcolor{ccwcolor}{\rule{8pt}{8pt}} CCW (counter-clockwise)};

\end{tikzpicture}
\end{document}
}
    \end{column}
    \begin{column}{0.5\textwidth}
      \begin{table}
        \centering\footnotesize
        \begin{tabular}{lcccc}
          \hline
          \textbf{Motor} & \textbf{T} & \textbf{Roll} & \textbf{Pitch} & \textbf{Yaw} \\
          \hline
          M1 FR & + & + & $-$ & $-$ \\
          M2 RR & + & + & + & + \\
          M3 RL & + & $-$ & + & $-$ \\
          M4 FL & + & $-$ & $-$ & + \\
          \hline
        \end{tabular}
      \end{table}

      \vspace{0.5em}

      \begin{exampleblock}{\footnotesize 直感的に理解する}
        \scriptsize
        Pitch+ → 前傾 → 前(M1,M4)弱 + 後(M2,M3)強\\
        Roll+ → 右傾 → 左(M3,M4)弱 + 右(M1,M2)強
      \end{exampleblock}
    \end{column}
  \end{columns}
\end{frame}

% ------------------------------------------------------------------
\begin{frame}{オープンループ制御 / Open-Loop Control}
  \begin{center}
    \resizebox{0.75\textwidth}{!}{\documentclass[border=10pt]{standalone}
\usepackage{tikz}
\usetikzlibrary{arrows.meta, positioning, shapes.geometric}

\begin{document}

\definecolor{blockcolor}{RGB}{0, 120, 200}
\definecolor{arrowcolor}{RGB}{80, 80, 80}
\definecolor{notefg}{RGB}{180, 60, 20}
\begin{tikzpicture}[
    block/.style={
      rectangle, draw=blockcolor, line width=1.5pt,
      fill=blockcolor!8, rounded corners=3pt,
      minimum width=28mm, minimum height=14mm,
      font=\bfseries, text=blockcolor, align=center
    },
    arrow/.style={-{Stealth[length=3mm]}, line width=1.2pt, arrowcolor},
    lbl/.style={font=\small, above, text=arrowcolor},
  ]

  % Blocks
  \node[block] (stick) at (0, 0) {Controller\\Stick};
  \node[block, right=25mm of stick] (scale) {Scaling\\$\times$ gain};
  \node[block, right=25mm of scale] (motor) {Motor\\PWM};
  \node[block, right=25mm of motor] (drone) {StampFly\\Drone};

  % Arrows
  \draw[arrow] (stick) -- node[lbl] {$-1.0 \sim +1.0$} (scale);
  \draw[arrow] (scale) -- node[lbl] {$0.0 \sim 1.0$} (motor);
  \draw[arrow] (motor) -- node[lbl] {duty} (drone);

  % No feedback annotation
  \node[font=\large\bfseries, text=notefg, below=18mm of scale]
    (nofb) {No Feedback!};
  \draw[notefg, line width=1pt, dashed, -{Stealth[length=2mm]}]
    (nofb.east) to[out=0, in=-90] ++(3, 1);

  % Title
  \node[font=\Large\bfseries, above=12mm of scale]
    {Open-Loop Control};

\end{tikzpicture}
\end{document}
}
  \end{center}

  \vspace{0.5em}

  \begin{alertblock}{問題点}
    フィードバックがないため、外乱（風・重心ずれ）に対応できない！\\
    → Lesson 5--6 で PID 制御を導入して解決
  \end{alertblock}
\end{frame}

% ------------------------------------------------------------------
\begin{frame}[fragile]{実習: 手動ミキシング / Hands-on: Manual Mixing}
  \begin{lstlisting}[style=sfcpp, title={student.cpp}]
#include "workshop_api.hpp"
static uint32_t tick = 0;
void setup() { ws::print("L2: Open-Loop Control"); }
void loop_400Hz(float dt) {
    tick++;
    float t = ws::rc_throttle();
    float r = ws::rc_roll()  * 0.3f;
    float p = ws::rc_pitch() * 0.3f;
    float y = ws::rc_yaw()   * 0.3f;
    ws::motor_set_duty(1, t + r - p - y); // FR
    ws::motor_set_duty(2, t + r + p + y); // RR
    ws::motor_set_duty(3, t - r + p - y); // RL
    ws::motor_set_duty(4, t - r - p + y); // FL
    if (tick % 80 == 0)
        ws::print("T=%.2f R=%.2f P=%.2f Y=%.2f",
            t, r, p, y);
}
  \end{lstlisting}
\end{frame}

% ------------------------------------------------------------------
\begin{frame}{チェックポイント / Checkpoint}
  \begin{alertblock}{確認事項}
    \begin{itemize}
      \item[$\square$] コントローラとペアリングできた
      \item[$\square$] スロットルで全モータが均等に回る
      \item[$\square$] ピッチスティックで前後モータの回転差が出る
      \item[$\square$] ゲイン(0.3f)を変えて挙動の違いを確認した
    \end{itemize}
  \end{alertblock}

  \vspace{1em}

  \begin{block}{次のレッスン}
    Lesson 3: LED 制御 --- システム状態を LED で可視化
  \end{block}
\end{frame}

\end{document}
